\documentclass[10pt]{article}

\usepackage[a4paper,top=13mm,bottom=12mm,left=18mm,right=18mm]{geometry}
\usepackage[T1]{fontenc}
\usepackage[utf8]{inputenc}
\usepackage[english]{babel}
\usepackage{lmodern}
\usepackage{microtype}
\usepackage[table]{xcolor}
\usepackage{tabularx}
\usepackage{array}
\usepackage{booktabs}
\usepackage{ragged2e}
\usepackage{eso-pic}

\definecolor{wine}{HTML}{8A2636}
\definecolor{ink}{HTML}{25282D}
\definecolor{muted}{HTML}{777A80}
\definecolor{rulegray}{HTML}{C9C9C9}
\definecolor{phasefill}{HTML}{F3F3F3}

\color{ink}
\setlength{\parindent}{0pt}
\setlength{\parskip}{2.15pt}
\setlength{\tabcolsep}{3.5pt}
\renewcommand{\arraystretch}{1.20}
\pagestyle{empty}

\newcommand{\proposalheading}[1]{%
  \vspace{1.8mm}%
  {\color{wine}\sffamily\fontsize{9.8}{11.2}\selectfont\bfseries
  \MakeUppercase{#1}}\par
  \vspace{0.55mm}%
}

\newcolumntype{F}{>{\raggedright\arraybackslash\color{wine}\sffamily\bfseries\fontsize{7.8}{9}\selectfont}p{25mm}}
\newcolumntype{A}{>{\raggedright\arraybackslash\fontsize{7.8}{9}\selectfont}X}
\newcolumntype{E}{>{\raggedright\arraybackslash\fontsize{7.8}{9}\selectfont}p{50mm}}

\begin{document}
\AddToShipoutPictureBG*{%
  \AtPageUpperLeft{\raisebox{-\height}{\color{wine}\rule{\paperwidth}{4mm}}}%
}

\begin{minipage}[t][0.98\textheight][t]{\textwidth}
\vspace{3.5mm}

\begin{center}
  {\fontsize{24}{27}\selectfont University Course Timetabling\par}
  \vspace{2.0mm}
  {\color{muted}\sffamily\fontsize{8.8}{10.2}\selectfont\bfseries
  STATE OF THE ART AND REPRODUCIBILITY IN\par
  CURRICULUM-BASED COURSE TIMETABLING\par}
  \vspace{2.0mm}
  {\color{wine}\rule{36mm}{0.8pt}}
\end{center}

\proposalheading{Problem and objective}
{\fontsize{9.4}{11.2}\selectfont\justifying
University timetabling assigns lectures to periods and rooms under teacher, room, and curriculum constraints. Timetable quality also depends on room capacity, weekly distribution, curriculum compactness, and room stability. This study focuses on \textbf{Curriculum-Based Course Timetabling (CB-CTT)} and the public UD2 formulation from ITC-2007. Its aim is to review current methods and assess whether recent open-source results can be reproduced.\par}

\proposalheading{Scope of the review}
{\fontsize{9.2}{11}\selectfont\justifying
The review will cover CB-CTT research published from 2015 to September 2026, together with the foundational work behind the benchmark. Searches in Scopus, Web of Science, and IEEE Xplore will be extended through citation tracing. Studies will be compared by formulation, dataset, method, evidence, computational setting, and code availability. The review will cover exact and constraint methods, metaheuristics, hybrid approaches, and learning-assisted techniques. ITC-2019 will provide a more recent point of comparison because it includes student sectioning and richer institutional constraints.\par}

\vspace{0.7mm}
\noindent\begin{tabularx}{\linewidth}{@{}F A E@{}}
\toprule[0.65pt]
\rowcolor{phasefill}
\color{wine}\sffamily\bfseries FAMILY &
\color{wine}\sffamily\bfseries REVIEW COVERAGE &
\color{wine}\sffamily\bfseries RECENT DIRECTION\\
\midrule[0.45pt]
Exact and constraint & MIP, CP, CP-SAT, decomposition, bounds & Model-based construction and decomposition\\
\midrule[0.35pt]
Metaheuristics & SA, tabu search, VNS/VND, ALNS, hyper-heuristics & LectureKick Simulated Annealing [5]\\
\midrule[0.35pt]
Hybrid methods & Solver-assisted construction and local improvement & CP-SAT initialization followed by VND [6]\\
\midrule[0.35pt]
Learning-assisted & Feature-based tuning, algorithm selection, adaptive control & Instance-space and algorithm-selection studies [3]\\
\bottomrule[0.65pt]
\end{tabularx}

\proposalheading{Reproducibility study}
{\fontsize{9.2}{11}\selectfont\justifying
The practical work will reproduce selected experiments from Ceschia et al. [5]. Their public repository includes source code, benchmark instances, and the parameter configuration used in the study. We will compile and run the solver on the \textbf{19 non-redundant ITC-2007 UD2 instances}, compare objective values and running times with the published results, and document the execution environment and any discrepancies. Instance \texttt{comp11} is omitted because competitive methods consistently reach cost zero. Instance \texttt{comp15} duplicates \texttt{comp03} under UD2 [1].\par}

\proposalheading{Expected contribution and scope}
{\fontsize{9.2}{11}\selectfont\justifying
The final article will combine a critical review of CB-CTT with an assessment of the reproducibility of a recent open-source solver. The discussion will address the strength of the available evidence, reporting transparency, and the limits of applying ITC-2007 results to present-day university timetabling.\par}

\proposalheading{References}
{\fontsize{7.2}{8.3}\selectfont\raggedright
[1] Ceschia, Di Gaspero, and Schaerf, ``Educational timetabling: Problems, benchmarks, and state-of-the-art results,'' 2023. \quad
[2] Abdipoor et al., ``Metaheuristic approaches for university timetabling problems,'' 2023. \quad
[3] De la Rosa-Rivera et al., ``Algorithm selection for educational timetabling,'' 2021.\par
[4] Gu et al., ``From integer programming to machine learning: A survey on solving university timetabling problems,'' 2025. \quad
[5] Ceschia et al., ``Revisiting local search for curriculum-based course timetabling,'' 2026. \quad
[6] Nguyen, ``Adaptive penalty-guided VND with CP-SAT for curriculum-based course timetabling,'' 2026. \quad
[7] Di Gaspero, McCollum, and Schaerf, \textit{The Second International Timetabling Competition (ITC-2007): Curriculum-Based Course Timetabling}, 2007.}

\vfill
{\color{rulegray}\rule{\linewidth}{0.55pt}}\vspace{0.6mm}
{\color{muted}\sffamily\fontsize{7.1}{8.2}\selectfont
Methodology and Research Techniques \textbar{} Practice 1 \textbar{} 2026-2027
\hfill Research proposal}

\end{minipage}
\end{document}
